% ===== 5 实验与结果 =====
\section{实验设计与阶段性结果}
\label{sec:r-results}

\subsection{数据角色、证据协议与评价指标}

本文将实验分为三类证据。\textbf{功能演示}回答处理链能否从真实自然堆积图像稳定地产生实例、毫米几何、级配、形貌和异常输出；\textbf{软件一致性}通过合成单元测试检查数学定义和代码实现是否一致；\textbf{物理精度验证}则必须使用独立 ground truth，才能声明实例精度、筛分精度或形貌/异常准确率。三者不能互相替代，特别是没有机械筛分称量真值时，视觉级配曲线只能视为模型输出而非“准确率”。

表\ref{tab:r-evidence-roles}把当前已有数据与正式比赛仍需补充的数据分开。三张演示图用于确认完整数据链和展示场景差异；人工实例掩膜应覆盖普通区域与粘连/遮挡困难区域；机械筛分数据必须与成像使用同一物理批次；形貌与异常标签则用于附加任务。这样的拆分防止最常见的验证错误：用模型自己的另一个输出作为“真值”，或把同一物理批次的多张照片随机拆到训练和测试两侧。

\begin{table}[t]
\centering
\caption{评委版实验中不同数据的角色与当前状态。}
\label{tab:r-evidence-roles}
\footnotesize
\renewcommand{\arraystretch}{1.12}
\begin{tabular}{>{\raggedright\arraybackslash}p{3.0cm}>{\raggedright\arraybackslash}p{4.3cm}>{\raggedright\arraybackslash}p{4.2cm}>{\raggedright\arraybackslash}p{2.0cm}}
\toprule
数据 & 用途 & 关键输出/指标 & 状态 \\
\midrule
3 张自然堆积 demo & 端到端功能、行为性分析、可视化 & 实例数、$D_p$、粒级、形貌/异常 & 已有 \\
合成单元测试 & 数学定义与代码一致性 & 27 项测试是否通过 & 已有 \\
人工实例 mask & 上游实例精度与 merge/split & P/R/F1、IoU、merge/split & 待补 \\
配对机械筛分批次 & $\boldsymbol{\theta},\gamma$ 校准与独立筛分评价 & $D_p$ 误差、逐筛级误差、曲线误差 & 关键缺口 \\
人工形貌/异常标签 & 附加能力外部有效性 & Macro-F1、P/R/F1 & 待补 \\
完整现场演练 & 30~min 工程约束 & 端到端 wall-clock time、失败恢复 & 待实测 \\
\bottomrule
\end{tabular}
\end{table}

最终正式实验应围绕同一物理批次建立配对证据：相机分支产生视觉预测，机械筛分分支产生逐筛质量真值。校准批次用于拟合 $\boldsymbol{\theta},\gamma$；独立测试批次在参数冻结后计算
\begin{equation}
 E_{D_p}=\left|\widehat D_p-D_p^{\mathrm{GT}}\right|,
 \qquad
 \mathrm{MAE}_{\mathrm{sieve}}=\frac{1}{K}\sum_{k=1}^{K}\left|\widehat P_k-P_k^{\mathrm{GT}}\right|.
\label{eq:r-metrics}
\end{equation}
还应同时报告相对 $D_p$ 误差和累计通过率曲线，以避免单一指标掩盖局部筛级偏差。实例识别使用人工实例 mask 报告 Precision/Recall/F1、IoU 及 merge/split；二维形貌和异常分别使用人工标签报告 Macro-F1 或 Precision/Recall。工程性能按完整 SOP 的 wall-clock time 评价，而不是只统计神经网络前向时间。

实验采样还需要覆盖真正可能改变误差结构的场景变量，而不是只增加相似照片数量。至少应包含偏细、连续和偏粗级配；低/中/高接触程度；不同位置与轻度光照变化；以及人为加入的超限块和泥团候选。若最终机械筛分批次数量有限，宁可保证物理批次之间的级配差异和严格独立测试，也不应通过同一批材料拍很多相似图片制造虚假的大样本量。

\subsection{真实图像功能演示与场景差异}

当前仓库包含 \texttt{agg\_001}、\texttt{agg\_005} 和 \texttt{agg\_029} 三张自然堆积演示图像，分辨率均为 3024$\times$4032，使用已写入产物的尺度 0.208~mm/px。三张图在当前基线 $\boldsymbol{\theta}=(1,0,0)$、$\gamma=3$ 下共检测 3858 个实例，并生成逐颗粒几何表、筛分等效粒径、质量代理级配、形貌/异常标签及可视化。表\ref{tab:r-demo-main}汇总正常实例口径下的当前输出。由于参数尚未用本赛题配对机械筛分批次校准，这些 $D_p$ 只能解释为\textbf{当前视觉质量代理预测}。

\begin{table}[t]
\centering
\caption{三张自然堆积演示样本的当前视觉质量代理输出。}
\label{tab:r-demo-main}
\footnotesize
\renewcommand{\arraystretch}{1.12}
\begin{tabular}{lrrrrll}
\toprule
样本 & 实例数 & 正常实例 & $D_{10}$ & $D_{50}$ & $D_{90}$ & 主粒级 \\
 & & & \multicolumn{3}{c}{/ mm} & \\
\midrule
\texttt{agg\_001} & 1924 & 1193 & 11.1 & 23.6 & 30.6 & 25--31.5~mm（34.2\%） \\
\texttt{agg\_005} & 968 & 703 & 15.6 & 27.5 & 38.0 & 25--31.5~mm（29.7\%） \\
\texttt{agg\_029} & 966 & 520 & 6.7 & 14.6 & 30.6 & 5--10~mm（32.7\%） \\
\bottomrule
\end{tabular}
\end{table}

三张场景得到明显不同的预测结构：\texttt{agg\_005} 粗粒尾部更强，\texttt{agg\_029} 则明显偏细，$D_{50}$ 排序为 \texttt{agg\_029}<\texttt{agg\_001}<\texttt{agg\_005}。这支持“同一处理链能将不同真实场景转换为非退化的场景级级配描述”，但不能说明哪条曲线最接近机械筛分。图\ref{fig:r-three-gsd}进一步把三个场景的数量/质量代理输出并列展示，可以看到差异不是由一个汇总数字产生，而是在完整分布形状上存在。

\begin{figure}[t]
\centering
\begin{subfigure}{0.32\linewidth}
\centering
\includegraphics[width=\linewidth]{agg_001_IG2_full_set_cp_SAM_pred_grains_re_scaled_gsd_comparison.png}
\caption{\texttt{agg\_001}}
\end{subfigure}\hfill
\begin{subfigure}{0.32\linewidth}
\centering
\includegraphics[width=\linewidth]{agg_005_IG2_full_set_cp_SAM_pred_grains_re_scaled_gsd_comparison.png}
\caption{\texttt{agg\_005}}
\end{subfigure}\hfill
\begin{subfigure}{0.32\linewidth}
\centering
\includegraphics[width=\linewidth]{agg_029_IG2_full_set_cp_SAM_pred_grains_re_scaled_gsd_comparison.png}
\caption{\texttt{agg\_029}}
\end{subfigure}
\caption{三张演示样本的数量型与质量代理型粒径分布。图中没有机械筛分真值，因此用于比较场景行为和统计口径，不作为物理准确率曲线。}
\label{fig:r-three-gsd}
\end{figure}

从逐筛质量代理比例也能看出场景结构差异。\texttt{agg\_001} 的 20--31.5~mm 合计约 55.3\%，\texttt{agg\_005} 的 25--40~mm 约 55.6\%，而 \texttt{agg\_029} 的 5--16~mm 约 51.9\%。这些比例来自同一计算规则，因此能够用于检查实现是否对不同输入给出一致口径；但它们只有在机械筛分配对后才可转化为“预测误差”。

\subsection{数量统计与质量代理：核心行为性证据}

当前最稳定且最重要的观察是：\textbf{数量型粒径分布与质量代理型粒径分布存在显著差异}。表\ref{tab:r-number-mass}给出异常处理前、基于全部检测实例的分位数。三张图的数量型 $D_{50}$ 仅为 5.2--7.6~mm，而采用 $w=d^3$ 后移动到 14.1--27.5~mm；方向在三个场景中一致。

\begin{table}[t]
\centering
\caption{数量型与质量代理型分位数比较。该结果证明统计口径影响显著，不证明 $\gamma=3$ 已获得物理校准。}
\label{tab:r-number-mass}
\footnotesize
\renewcommand{\arraystretch}{1.12}
\begin{tabular}{lrrrrrr}
\toprule
& \multicolumn{3}{c}{数量型 / mm} & \multicolumn{3}{c}{质量代理 / mm} \\
\cmidrule(lr){2-4}\cmidrule(lr){5-7}
样本 & $D_{10}$ & $D_{50}$ & $D_{90}$ & $D_{10}$ & $D_{50}$ & $D_{90}$ \\
\midrule
\texttt{agg\_001} & 3.1 & 6.0 & 17.3 & 10.6 & 23.5 & 30.6 \\
\texttt{agg\_005} & 3.3 & 7.6 & 24.5 & 15.4 & 27.5 & 38.0 \\
\texttt{agg\_029} & 3.0 & 5.2 & 9.2 & 5.9 & 14.1 & 30.6 \\
\bottomrule
\end{tabular}
\end{table}

该变化可由式\eqref{eq:r-mass}直接解释：自然堆积图中小实例数量很多，但其单颗粒质量贡献较低；少量粗颗粒虽然在个数中占比小，却在幂次权重下显著影响累计质量。以 \texttt{agg\_001} 为例，当前 $<5$~mm 实例占检测数量的 38.0\%，剔除这些实例后质量代理 $D_{50}$ 仅由 23.50~mm 变为 23.56~mm，说明“大量小实例”不必然意味着“大量质量”。相反，错误 merge 产生的少量粗目标会被 $d^{\gamma}$ 放大，因此独立实例边界对高分位质量级配尤为关键。

\texttt{agg\_029} 提供了另一个有用现象：其正常实例仅占检测数的 53.8\%，过小候选数量最多，但正常口径下质量代理 $D_{90}$ 仍达到 30.6~mm。换言之，数量端“看起来很细”和质量端仍存在粗颗粒贡献可以同时成立。这正是机械筛分按质量计量与机器视觉按目标计数之间的统计差异，也是本文必须显式建立质量映射而非直接报告实例直方图的原因。

需要区分两个结论：现有结果足以支持“不能忽略质量口径”；但 $\gamma$ 越大本来就会推动分布向粗粒移动，因此不能仅凭曲线看起来更合理就选定 $\gamma=3$。只有当某个 $\gamma$ 在独立机械筛分测试上降低式\eqref{eq:r-metrics}的误差时，才能把它称为更好的物理质量模型。

\subsection{投影形貌与异常候选}

当前几何规则在三张图上均产生了非退化的投影形貌分布，具体比例见表\ref{tab:r-shape-results}。普通类占 55.0\%--64.1\%，针片状候选约 14.0\%--20.9\%；\texttt{agg\_005} 的针片状候选比例最高，而 \texttt{agg\_029} 的圆形比例相对较高。这些差异说明 AR、圆度和凸度能够把场景中的二维轮廓组织成可检查统计，但没有人工/量规 ground truth 时不能进一步解释为材料真实针片状率。

\begin{table}[t]
\centering
\caption{三张演示样本的投影形貌规则输出（占全部检测实例比例，\%）。}
\label{tab:r-shape-results}
\footnotesize
\begin{tabular}{lrrrr}
\toprule
样本 & 针片状候选 & 圆形 & 棱角状 & 普通 \\
\midrule
\texttt{agg\_001} & 16.3 & 8.0 & 16.4 & 59.3 \\
\texttt{agg\_005} & 20.9 & 6.6 & 17.6 & 55.0 \\
\texttt{agg\_029} & 14.0 & 11.0 & 11.0 & 64.1 \\
\bottomrule
\end{tabular}
\end{table}

图\ref{fig:r-shape-montage}以 \texttt{agg\_001} 为例展示四类高亮结果。图的价值在于检查规则是否定位到了与定义一致的轮廓类型，而不是用视觉印象代替准确率。例如，针片状候选应主要集中于低 AR 的细长投影；圆形应同时具有较高圆度和凸度。若高亮区域明显不符合这些定义，优先应检查实例 mask 和阈值实现，而不是直接解读最终比例。

\begin{figure}[t]
\centering
\begin{subfigure}{0.24\linewidth}\centering
\includegraphics[width=\linewidth]{agg_001_IG2_full_set_cp_SAM_pred_grains_re_scaled_shape_needle_flaky.png}
\caption{针片状候选}
\end{subfigure}\hfill
\begin{subfigure}{0.24\linewidth}\centering
\includegraphics[width=\linewidth]{agg_001_IG2_full_set_cp_SAM_pred_grains_re_scaled_shape_round.png}
\caption{圆形}
\end{subfigure}\hfill
\begin{subfigure}{0.24\linewidth}\centering
\includegraphics[width=\linewidth]{agg_001_IG2_full_set_cp_SAM_pred_grains_re_scaled_shape_angular.png}
\caption{棱角状}
\end{subfigure}\hfill
\begin{subfigure}{0.24\linewidth}\centering
\includegraphics[width=\linewidth]{agg_001_IG2_full_set_cp_SAM_pred_grains_re_scaled_shape_regular.png}
\caption{普通}
\end{subfigure}
\caption{\texttt{agg\_001} 的二维投影形貌高亮。类别由几何阈值产生，图像用于规则可解释性检查，不替代专家形貌 ground truth。}
\label{fig:r-shape-montage}
\end{figure}

异常分支的当前输出见表\ref{tab:r-anomaly-results}。三张图均未触发 $>50$~mm oversized 或疑似泥团规则，主要触发的是 $<5$~mm 小目标。该结果不能说明系统对 oversized 或泥团“零误报且完全正确”，更准确的结论是：\textbf{当前三张演示数据不包含足以检验这两个分支的阳性证据。} 正式测试应主动构造正常骨料、已知 $>50$~mm 大块和泥团/非骨料异物混合样本，分别计算 Precision/Recall，而不是等待随机场景恰好出现异常。

\begin{table}[t]
\centering
\caption{演示样本的异常规则输出（实例数）。}
\label{tab:r-anomaly-results}
\footnotesize
\begin{tabular}{lrrrr}
\toprule
样本 & 正常 & $>50$~mm & $<5$~mm & 疑似泥团 \\
\midrule
\texttt{agg\_001} & 1193 & 0 & 731 & 0 \\
\texttt{agg\_005} & 703 & 0 & 265 & 0 \\
\texttt{agg\_029} & 520 & 0 & 446 & 0 \\
\bottomrule
\end{tabular}
\end{table}

\subsection{软件一致性与工程时序}

仓库为筛分等效粒径、质量权重、加权分位数、形貌规则、异常优先级和报告生成提供了 27 项合成单元测试，当前记录全部通过。表\ref{tab:r-software-tests}说明这些测试能够证明和不能证明什么。特别是，合成数据可以构造已知 $\boldsymbol{\theta},\gamma$ 并检查优化器是否恢复接近设定值，这证明参数拟合路径在数学上可运行；它并不证明真实骨料确实服从该模型。

\begin{table}[t]
\centering
\caption{竞赛层软件测试的证据范围。}
\label{tab:r-software-tests}
\footnotesize
\renewcommand{\arraystretch}{1.12}
\begin{tabular}{>{\raggedright\arraybackslash}p{3.2cm}>{\raggedright\arraybackslash}p{5.0cm}>{\raggedright\arraybackslash}p{5.2cm}}
\toprule
模块 & 测试关注点 & 不能替代的现实验证 \\
\midrule
筛分等效与校准 & 粒径公式、质量权重、分位数、合成参数恢复 & 真实二维粒径与机械筛孔的对应关系 \\
形貌 & AR/$C$/$S$ 阈值与优先级逻辑 & 专家/量规形貌准确率 \\
异常 & oversized/undersized/泥团候选规则边界 & 真实异常 Precision/Recall \\
报告 & JSON/表格字段、正常口径与异常剔除一致性 & 现场输入质量和外部物理准确率 \\
\bottomrule
\end{tabular}
\end{table}

历史运行记录约为实例分割 17~s、几何测量 9~s、筛分统计 $<1$~s，即算法阶段约 27~s/张。计算本身因此不是 30~min 约束下的主要瓶颈；更可能占用现场时间的是设备就位、尺度确认、采集质检、模型首次加载、结果目检和一次失败恢复。基于这一判断，第\ref{sec:r-deploy}节把 30~min 预算分配给完整操作链，而不是将“27 秒”直接写成系统总耗时。

\subsection{最终机械筛分验证与最小消融矩阵}

当前最关键的实验缺口是：三张演示样本没有配对提交逐筛机械称量数据。因此本版本不填写机械筛分准确率，也不把未校准 $D_p$ 作为竞赛精度结论。正式测试应至少建立若干配对批次，其中一部分用于拟合参数，另一部分完全隔离为测试集。独立测试中至少报告：三个 $D_p$ 的绝对/相对误差、逐筛级质量占比或累计通过率误差，以及各批次均值与离散程度。表\ref{tab:r-physical-template}给出最终报告应替换的结果位置；当前保持“待测”比填入视觉自比较数字更符合证据逻辑。

\begin{table}[t]
\centering
\caption{独立机械筛分测试的最终结果模板。获得同批逐筛称量后直接回填。}
\label{tab:r-physical-template}
\footnotesize
\renewcommand{\arraystretch}{1.12}
\begin{tabular}{lrrrrr}
\toprule
测试批次 & $|\Delta D_{10}|$ & $|\Delta D_{50}|$ & $|\Delta D_{90}|$ & 筛级 MAE & 曲线/累计误差 \\
 & \multicolumn{3}{c}{/ mm} & / percentage point & \\
\midrule
Test-1 & 待测 & 待测 & 待测 & 待测 & 待测 \\
Test-2 & 待测 & 待测 & 待测 & 待测 & 待测 \\
Test-3 & 待测 & 待测 & 待测 & 待测 & 待测 \\
\midrule
平均 & 待测 & 待测 & 待测 & 待测 & 待测 \\
\bottomrule
\end{tabular}
\end{table}

同时需要用同一独立测试集完成表\ref{tab:r-ablation-matrix}的最小消融。实验顺序应回答四个因果问题：短轴是否比面积等效直径更接近筛分；二者组合是否进一步降低误差；质量加权是否比数量等权更接近称重结果；机械筛分校准是否能在未见批次上继续降低误差。异常剔除则用于检查极端误检是否通过 $d^\gamma$ 放大并污染最终级配。

\begin{table}[t]
\centering
\caption{正式物理验证的最小消融矩阵。所有行必须在同一独立测试批次上比较。}
\label{tab:r-ablation-matrix}
\footnotesize
\renewcommand{\arraystretch}{1.1}
\begin{tabular}{>{\centering\arraybackslash}p{0.8cm}>{\raggedright\arraybackslash}p{4.0cm}>{\raggedright\arraybackslash}p{4.6cm}>{\raggedright\arraybackslash}p{4.0cm}}
\toprule
编号 & 设置 & 回答的问题 & 主要比较指标 \\
\midrule
A & $d=b$，数量等权 & 最简视觉粒径/计数基线有多大语义差距？ & $D_p$、筛级 MAE \\
B & $d=b$，固定 $\gamma=3$ & 仅引入质量口径是否改善？ & B vs. A \\
C & $d=d_{\mathrm{eq}}$，同一 $\gamma$ & 面积尺度还是短轴更接近筛孔？ & C vs. B \\
D & 式\eqref{eq:r-sieve-d}，固定 $\gamma$ & 低维几何融合是否有效？ & D vs. B/C \\
E & 式\eqref{eq:r-sieve-d}，校准 $\gamma$ & 质量指数校准是否具有独立收益？ & E vs. D \\
F & E + 异常/质量控制 & 极端实例过滤能否降低粗端误差？ & F vs. E；失败案例 \\
\bottomrule
\end{tabular}
\end{table}

真正有解释力的结论应表述为“某设计使机械筛分独立测试误差降低多少”，而不是“参数变化后曲线移动了多少”。若 E 相比 D 没有稳定收益，则应保留更简单的固定质量模型；若 D 相比 $b$ 基线没有改善，则不应为了方法复杂度保留额外参数。换言之，最终系统结构也应由物理验证裁剪，而不是默认所有提出的模块都必须保留。在获得这些 ground truth 之前，本文明确保留证据边界，比给出未经验证的高准确率更能反映系统当前真实成熟度。
